% problem-000249 10 富勒烯材料与磁共振
\section*{10. 富勒烯材料与磁共振}
长期以来，人们认识到碳只有三种同素异形体，分别是金刚石、石墨与无定形碳。1985 年，Kroto、Curl 和 Smalley 在解析激光蒸发石墨获得的碳团簇质谱时，发现了碳的第四种同素异形体 $\mathbf { C } _ { 6 0 } .$ 。富勒烯的发现极大地推动了碳化学的发展，促进了与其相关的材料科学、纳米科学和电子学的发展。

\subsection*{10-1}
已知 $C _ { 6 0 }$ 的晶胞为面心立方结构，其晶胞参数 $a = 1 . 4 1 3 \mathrm { n m }$ 。比较 $\mathbf { C } _ { 6 0 }$ 、金刚石、石墨三者密度的大小。

\subsection*{10-2}
N可以与 $C _ { 6 0 }$ 分子结合，形成 ${ { C } _ { 6 0 } } \mathrm { { N } }$ 分子。研究发现 ${ \cal C } _ { 6 0 } \mathsf { N }$ 与 ${ C } _ { 6 0 }$ 的化学反应性几乎完全相同。

\subsection*{10-2}
-1不同自旋量子数的电子(S)或原子核(I)的磁能级在外加磁场时会发生分裂，即塞曼分裂。塞曼分裂能 $\Delta E$ 可通过旋磁比 $\gamma$ 与外加磁场B的乘积来计算。常见同位素基本信息如表10.1所示。写出N的两种同位素核自旋量子数（1）可能的取值$( m _ { I } )$ ，计算这两种同位素的不同 $m _ { I }$ 在0.3000T磁场下的能量差（用频率表示），并画出二者能级结构的示意图且注明能级间的能量差。

表 10.1 常见原子的同位素信息

\textit{（原卷此处含表格/仪器试剂表，详见 PDF 或 MD 源文件。）}

\subsection*{10-2}
-2电子顺磁共振（EPR）与核磁共振技术（NMR）相似，可用于研究电子自旋性质。当微波频率与电子的塞曼分裂能相同时，EPR的谱图会显示一个信号峰。与NMR 类似，磁性原子核 $( I > 0 )$ 同样会与电子产生耦合，使EPR中的电子自旋的谱峰进一步分裂为 $( 2 I + 1 )$ 重峰，图10.1是电子自旋磁能级的能量分裂示意图。

（图：）
图10.1 EPR 中电子自旋磁能级的能量分裂示意图
第15页 共17页

使用微波频率 9.834 GHz 测试得到的 ${ \cal C } _ { 6 0 } \Nu$ 分子EPR谱如图10.2。判断 $C _ { 6 0 } \mathrm { N }$ 分子中未成对电子的位置。

（图：）
图 10.2 ${ C _ { 6 0 } } ^ { 1 4 } \mathrm { N }$ 与 $\mathrm { C } _ { 6 0 } \mathrm { { } ^ { 1 5 } N }$ 的 EPR 谱图

\subsection*{10-2}
-3 画出 ${ \sf C } _ { 6 0 } { \sf N }$ 的分子结构示意图（用圆表示 $\mathbf { C } _ { 6 0 } )$ 及其中N的电子排布图，并标注电子自旋方向。

\subsection*{10-3}
富勒烯材料是一种非常有潜力的量子信息材料。在量子信息技术中，人们希望能对样品进行有效的初始化与操控。

布居数是能级上分布的微观粒子数目：初始化是使粒子处在某一能级上的布居数比例超过99%，经过初始化后的系统状态是纯态：操控是指将两个相邻能级的布居数按照一定比例 $r \left( 0 \sim 1 \right)$ 进行交换，使系统按期望演化。使用符号 $U ( m , m + 1 , r )$ 描述一次操控。设计合适的操控序列，可实现复杂的量子信息处理过程。一个典型的操控 $U ( m , m + 1 , 0 . 2 5 )$ 的示例如图10.3所示。

（图：）
图 10.3 量子信息处理过程中，对 m 与 m+1 能级操控的示例

\subsection*{10-3}
-1 计算298.15 K下 ${ \cal C } _ { 6 0 } \mathbf { N }$ 分子中未成对电子在外加350.9mT磁场时各个磁能级的布居数比例。（忽略核自旋的影响，结果保留4位有效数字）

\subsection*{10-3}
-2量子信息处理过程中，各个能级布居数相同的部分不会对结果产生影响，因

第38届中国化学奥林匹克（决赛）理论第一场试题 2024年10月 26日 广州

而相同的部分可整体抽提出并忽略。如果各个能级的布居数在抽提出相同数量时，剩下的部分是一个纯态，则称该体系处于赝纯态。设计一种操控序列，将10-3-1中的热平衡态转变成雁纯态。(忽略核自旋的影响)

\subsection*{10-4}
液氨的正常沸点是 4.22 K，液氮的汽化热是 0.0829 kJ mol。

\subsection*{104-1}
低温是量子信息处理的必备条件，目前获取超低温的常用方法是通过液氮冷却。计算在4.22K下需要外加多大的磁场使 ${ \mathsf { C } } _ { 6 0 } \mathbf { N }$ 的电子自旋实现初始化。（忽略核自旋的影响，精确到0.1T)

\subsection*{10-4}
-2实验室可利用低压环境降低液氦沸点，获取低于4.22K的温度。计算需要多大的压力可以使体系温度降低至2.00K？(结果保留3位有效数字)

\subsection*{10-4}
-3磁致冷是一种非常有希望实现无氮超低温环境的方式。通过一定的热力学循环过程，可将低温区域降低至0.1 K量级。镓酸钆 $\mathbf { G } \mathbf { d } _ { 3 } \mathbf { G } \mathbf { a } _ { 5 } \mathbf { O } _ { 1 2 }$ (GGG，密度 7.09 g$\mathsf { c m } ^ { - 3 } )$ 是目前常用的磁致冷材料。200mK、5T时，磁致冷材料的电子自旋可几乎完全处在最低能态。

计算 GGG 与 ${ \cal C } _ { 6 0 } \mathbf { N }$ (忽略核自旋的影响）材料从200 mK升温至 298.15 K过程中的摩尔磁熵变、单位质量(g)磁熵变与单位体积(cm³)磁熵变。(结果保留3位有效数字，忽略温度变化对体积的影响，可做其他合理近似）

\subsection*{10-4}
-4 磁致冷过程的热力学循环如图10.4所示。

（图：）
图10.4 磁致冷过程的热力学循环示意图

该循环的4个过程均可逆，过程中均无体积变化。判断4个过程中的熵变(ΔS)吸放热（O）、磁化强度变化（ΔM)。

\textit{（原卷此处含表格/仪器试剂表，详见 PDF 或 MD 源文件。）}

(填“正”、“负”、“0”或“无法判断”)

第17页 共17页

\subsection*{关于自主选拔在线平台}

自主选拔在线平台（原自主招生在线）创办于 2014 年，历史可追溯至 2008 年，总部位于北京，是国内极具影响力的拔尖创新人才培养咨询服务平台。涵盖清北升学、学科竞赛、强基计划、综合评价、英才少年班、港澳升学、三位一体、研学营地、初高中生涯规划、志愿填报、大学保研留学等业务。

自主选拔在线旗下拥有微信公众平台、网站门户、视频号、小红书、微博、社群、抖音等矩阵生态平台。平台活跃用户达 500万+量级，网站门户年度访问量超 1亿+。用户群体涵盖北京、广东、江苏、浙江、山东、上海、安徽、四川等 31 省市（自治区），全国超 95%以上的重点中学校长老师、家长考生及知名高KII校招办老师都在关注，在行业内影响力首屈一指。

自主选拔在线平台一直秉承“深度、专业、有态度”的创办理念，不断探索“K12 教育+互联网+大数据”的运营模式，尝试基于 AI大数据为广大中学和家长提供学科竞赛、强基计划、综合评价、港澳升学、工作。

平台自创办以来，为众多重点大学发现和推荐优秀生源，和全国数百所重点中学达成深度战略合作。累计举办大学招办宣讲会、公益升学讲座数千场，累计帮助百余万考生成功考入清北复交浙等重点大学。在家长考生、中学和社会各界具有很好的品牌口碑，曾荣获央广网“年度口碑影响力在线教育品牌”。

未来，自主选拔在线将立足于新高考改革，全面整合高校、中学及教科研机构等各方面资源，依托在线教育模式，打造更加全面、专业的拔尖人才贯通培养升学咨询服务平台。

强基在线

（图：）
自主选拔在线
名校综合评价
学科竞赛派
京考一点通
学在港澳
广东家长圈
江苏学生圈
浙考家长帮
齐鲁家长圈
太星大学生

咨询热线：010-5601
